\chapter{From Signal to Forecast}\label{ch:rs:from-signal-to-forecast}

A stock's composite score is $+2$. Is its expected return over the next month two basis points or two hundred? The composite of
chapter~14 ranks stocks; a portfolio optimiser, a risk budget and a trader need expected returns in units of return, per stock
and per horizon, and a manager needs to know what information ratio the forecasts can deliver. This chapter turns scores into
expected returns (calibration), states the law that links a forecast's quality to a portfolio's information ratio, and measures on a
planted forecast where the law's promise is lost. The answer to the opening question, for a score with an IC of 0.05 on a stock whose
monthly residual volatility is 7.7\%, is 77 basis points.

\section{Calibrating a score}

\begin{definition}[Forecast calibration, alpha scaling rule]\label{def:rs:from-signal-to-forecast:calibration}
\emph{Forecast calibration}\index{forecast calibration} is the mapping of a score into an expected return such that, among the names
given any forecast, the average realised return equals the forecast. The \emph{alpha scaling rule}\index{alpha scaling rule} sets the
expected \omterm{def:rs:anatomy-of-a-predictor:residual}{residual return} of a name to $\alpha_i = \mathrm{IC}\cdot\omega_i\cdot z_i$, with $z_i$ the standardised score and $\omega_i$ the
name's residual volatility over the horizon.
\end{definition}

\begin{proposition}[The scaling rule is a regression]\label{prop:rs:from-signal-to-forecast:scaling}
If the standardised score $z$ and the volatility-scaled \omterm{def:rs:anatomy-of-a-predictor:residual}{residual return} $r/\omega$ are jointly normal with correlation $\mathrm{IC}$ and
unit variances, the expected return given the score is $\mathrm{E}[r \mid z] = \mathrm{IC}\cdot\omega\cdot z$.
\end{proposition}

\begin{proof}
For jointly normal standardised variables, $\mathrm{E}[r/\omega \mid z] = \mathrm{IC}\cdot z$; multiply by $\omega$.
\end{proof}

The rule makes three assumptions visible: that the score is standardised, that the IC is the same for all names (so that riskier names
get proportionally larger forecasts), and that the relation is linear. Each can be checked. The chapter's planted forecast
(\texttt{rs\_forecast}) adds to the monthly market-adjusted returns of the 704 names that \texttt{firm.synthmkt} lists for ten years a true
expected return $0.05\,\omega_i s_i$, with $s_i$ a standard normal score redrawn each month and $\omega_i$ the name's residual volatility
(median 7.7\% a month). Estimated on the first five years, the IC is 0.055. On the last five, the Mincer--Zarnowitz regression of realised
returns on the scaled forecasts (Book~4, chapter~18) has a slope of 0.92, close to the 1 of a calibrated forecast, and an $R^2$ of
0.24\%: calibrated forecasts explain almost none of the variance of returns, and that is what an IC of 0.05 means. The same regression on the
raw score has a slope of 0.0039, the IC times a typical volatility: the score was never in units of return.

\begin{definition}[Calibration curve, isotonic regression]\label{def:rs:from-signal-to-forecast:curve}
A \emph{calibration curve}\index{calibration curve} plots the mean realised return of names grouped by forecast (in bins) against the mean
forecast. \emph{Isotonic regression}\index{isotonic regression} fits the best non-decreasing function of the score to realised returns by least
squares, by pooling adjacent groups that violate the order (Ayer, Brunk, Ewing, Reid and Silverman, 1955).
\end{definition}

Binned by score decile, the realised volatility-scaled returns follow the scaling rule's line, with the noise of 4\,224 observations a decile
(\cref{fig:rs:from-signal-to-forecast:calibration}). \omterm{def:rs:from-signal-to-forecast:curve}{Isotonic regression} gives a calibration without the linear assumption and pays for it in
steps: flat stretches where the data could not order neighbouring scores, and larger values in the tails, which are fitted on few names.
Between the two, a linear calibration with an estimated slope is the default; the curve is the check.

\begin{omfigure}
\begin{tikzpicture}
\begin{axis}[width=10cm,height=6cm,xlabel={score $z$},ylabel={realised return / residual volatility},tick label style={font=\small},
  label style={font=\small},xmin=-2.6,xmax=2.6,ymin=-0.22,ymax=0.14,grid=major,grid style={gray!25},
  yticklabel style={/pgf/number format/fixed,/pgf/number format/precision=2},scaled y ticks=false,legend pos=north west,
  legend style={font=\footnotesize},table/col sep=comma]
\addplot[omThm,thick,dashed] table[x=score,y=predicted]{figdata/research/15-from-signal-to-forecast/calibration.csv};
\addplot[only marks,mark=*,omDef,mark size=2pt] table[x=score,y=realised]{figdata/research/15-from-signal-to-forecast/calibration.csv};
\addplot[black!60,thick,const plot] table[x=score,y=fit]{figdata/research/15-from-signal-to-forecast/isotonic.csv};
\legend{scaling rule $\widehat{\mathrm{IC}}\cdot z$,score deciles (realised),isotonic regression}
\end{axis}
\end{tikzpicture}

\omcaption{Calibration of the planted score on the last five years: mean realised return over residual volatility by score decile against the
scaling rule with the IC estimated on the first five (0.055), and the isotonic fit. Data: \texttt{rs\_forecast} on
\texttt{firm.synthmkt}.}\label{fig:rs:from-signal-to-forecast:calibration}
\end{omfigure}

\section{Scaling across assets and horizons}

The rule scales the forecast by each name's volatility, so a forecast is comparable across assets only through its IC. A signal that is
equally informative about a utility and a biotechnology stock forecasts a larger return for the biotechnology stock, because its returns are
larger; a forecast that gave both the same expected return would imply a much higher IC for the utility. Across horizons the scaling follows
the \omterm{def:rs:half-life-decay-and-stability:decay}{IC decay curve} of chapter~13: the expected return over $h$ periods is $\omega\,z$ times the sum of the single-period ICs up to $h$ (for
volatility per period $\omega$), not the one-period forecast times $h$. A signal whose information lasts a day forecasts the same expected
return over a month as over a day. Across asset classes, the standardisation must be within each class, and the IC estimated per class; a
blend of an equity and a bond forecast adds expected returns, never scores.

\section{The fundamental law of active management}

\begin{definition}[Breadth, fundamental law of active management]\label{def:rs:from-signal-to-forecast:law}
The \emph{breadth}\index{breadth} of a strategy is the number of independent forecasts it acts on per year. The \emph{fundamental law of active
management}\index{fundamental law of active management} (Grinold, 1989) states that the information ratio of the optimal portfolio built on
forecasts with \omterm{def:rs:anatomy-of-a-predictor:ic}{information coefficient} $\mathrm{IC}$ and breadth $\mathrm{BR}$ is $\mathrm{IR} \approx \mathrm{IC}\sqrt{\mathrm{BR}}$.
\end{definition}

With 704 names forecast each month and an IC of 0.05, the law promises $0.05\sqrt{12 \times 704} = 4.60$. The book that acts on the planted
forecast holds each name in proportion to $\alpha_i/\omega_i^2$ (the mean-variance weights for uncorrelated residuals), rebalanced monthly. When
the residuals are replaced by Gaussian noise of the same volatilities, independent across names, the book's realised information ratio over ten
years is 4.92, within its sampling error of the law (about 0.45 for ten years of monthly returns). The law holds where its assumptions do.

\section{Where the law breaks}

\begin{definition}[Effective breadth, transfer coefficient]\label{def:rs:from-signal-to-forecast:breadth}
The \emph{effective breadth}\index{effective breadth} of a set of forecasts is the number of independent forecasts that would give the same
information ratio; $n$ bets with average pairwise correlation $\rho$ are worth about $n/(1 + (n - 1)\rho)$. The \emph{transfer
coefficient}\index{transfer coefficient} of a portfolio is the cross-sectional correlation between its risk-adjusted active weights $w_i\omega_i$
and the risk-adjusted forecasts $\alpha_i/\omega_i$.
\end{definition}

\begin{proposition}[The generalised law]\label{prop:rs:from-signal-to-forecast:generalised}
Under the assumptions of the law, a portfolio whose risk-adjusted active weights have correlation $\mathrm{TC}$ with the risk-adjusted forecasts
has an expected information ratio of $\mathrm{TC}\cdot\mathrm{IC}\sqrt{\mathrm{BR}}$.
\end{proposition}

\begin{proof}
Write $x_i = w_i\omega_i$ and $u_i = r_i/\omega_i$, so that the active return is $\sum_i x_iu_i$ with $\mathrm{E}[u_i] = \mathrm{IC}\,z_i$ and independent
unit-variance noise. For a fixed $x$ with $\sum x_i^2 = 1$, the expected active return is $\mathrm{IC}\sum x_iz_i = \mathrm{IC}\sqrt n\,\mathrm{TC}$ (the
correlation of $x$ with $z$, both with norm fixed) and the risk is 1; over $\mathrm{BR}/n$ independent periods a year the ratio scales by
$\sqrt{\mathrm{BR}/n}$, giving $\mathrm{TC}\cdot\mathrm{IC}\sqrt{\mathrm{BR}}$.
\end{proof}

Clarke, de Silva and Thorley (2002) derived this generalised version for constrained portfolios, with an ex post decomposition of performance into
the success of the forecasts and the noise of the constraints. On the chapter's market the book loses its promise in three steps
(\cref{fig:rs:from-signal-to-forecast:shortfall}):

\textbf{The IC is noisy.} Qian and Hua reformulated the law with a random IC: the information ratio is the IC's mean over its standard
deviation (as summarised by Zhang, Wang and Cao, 2021). With the market's own residuals (industry and style factors, volatility clustering) the
monthly IC has a mean of 0.0505 and a standard deviation of 0.0424, against 0.0331 with Gaussian residuals; the ratio, annualised, is 4.13, and the
book's realised information ratio is 4.26. The factors make neighbouring bets move together, and the 704 names are worth
$(\mu_{\mathrm{IC}}/\sigma_{\mathrm{IC}})^2/\mathrm{IC}^2 = 569$ independent forecasts a month.

\textbf{The portfolio cannot follow the forecasts.} A long-only book around an equal-weight benchmark cannot hold a name below zero. Asked for a
monthly tracking error of 4\%, it clips 47\% of its intended holdings at zero, and its \omterm{def:rs:from-signal-to-forecast:breadth}{transfer coefficient} falls to 0.86. The generalised law then
predicts $0.857 \times 4.13 = 3.54$.

\textbf{The constraint adds noise.} The book's realised information ratio is 3.00: the part of the clipped book that is not proportional to the
forecasts is uncorrelated with them and adds risk without return. From 4.60 to 3.00, two thirds of the promise survives.

\begin{omfigure}
\begin{tikzpicture}
\begin{axis}[width=12cm,height=5.8cm,ybar,bar width=18pt,ylabel={information ratio (annual)},tick label style={font=\small},
  label style={font=\small},xtick={0,1,2,3,4,5},
  xticklabels={{law},{book,\\Gaussian},{IC mean\\over sd},{book,\\market},{$\times$ TC,\\long-only},{book,\\long-only}},
  xticklabel style={align=center,font=\footnotesize},ymin=0,ymax=5.5,enlarge x limits=0.1,grid=major,grid style={gray!25},
  nodes near coords,nodes near coords style={font=\footnotesize,/pgf/number format/fixed,/pgf/number format/fixed zerofill,
  /pgf/number format/precision=2},table/col sep=comma]
\addplot[fill=omDef!55,draw=omDef] table[x=k,y=ir]{figdata/research/15-from-signal-to-forecast/shortfall.csv};
\end{axis}
\end{tikzpicture}

\omcaption{The information ratio of the planted forecast (IC 0.05, 704 names, monthly): the law's promise, the book's with Gaussian residuals, the
IC's mean over its standard deviation and the book's with the market's residuals, the generalised law with the long-only \omterm{def:rs:from-signal-to-forecast:breadth}{transfer coefficient}, and
the long-only book. Data: \texttt{rs\_forecast.law}.}\label{fig:rs:from-signal-to-forecast:shortfall}
\end{omfigure}

The tracking error is the lever: the \omterm{def:rs:from-signal-to-forecast:breadth}{transfer coefficient} falls from 0.98 at a requested 0.25\% a month to 0.86 at 4\%, and the information ratio
from 3.91 to 3.00, while the share of clipped holdings rises from 13\% to 47\%.

\begin{center}
\small
\begin{tabular}{@{}lccccc@{}}
\toprule
requested tracking error (a month) & 0.25\% & 0.5\% & 1\% & 2\% & 4\%\\
\midrule
holdings clipped at zero & 13\% & 28\% & 38\% & 44\% & 47\%\\
\omterm{def:rs:from-signal-to-forecast:breadth}{transfer coefficient} & 0.98 & 0.93 & 0.89 & 0.87 & 0.86\\
information ratio & 3.91 & 3.54 & 3.25 & 3.09 & 3.00\\
\bottomrule
\end{tabular}
\end{center}

\section{Evaluating forecasts}

A forecast is evaluated as a forecast before it is evaluated as a strategy: its IC (mean, dispersion, and their ratio), its calibration (the
Mincer--Zarnowitz slope near 1, the \omterm{def:rs:from-signal-to-forecast:curve}{calibration curve}), its decay (chapter~13), its stability (chapter~13), and its incremental value over the
firm's other forecasts (chapter~12). The law then says what portfolio to expect, and the gap between the law and the backtest is not a mystery: it
is the IC's noise, the constraints and the costs (chapter~27), each measurable. Michaud, Esch and Michaud (2020) argued that applications of
Grinold's law to portfolio design, adding securities or factors or trading more often, are often unreliable because they ignore estimation error and
the constraints of practice; the decomposition above is the check that keeps a manager honest.

\section{Predictor cards}

\begin{predictorcard}{Calibrated composite}\label{pred:rs:from-signal-to-forecast:calibrated}
\sfield{Definition} $\alpha_i = \widehat{\mathrm{IC}}\,\omega_i z_i$, $z$ the composite of chapter~14 standardised by date, $\widehat{\mathrm{IC}}$ its trailing
mean IC, $\omega_i$ the risk model's residual volatility over the horizon.
\sfield{Inputs and timestamps} The composite and the risk model as of the decision date; the IC estimated on earlier dates only.
\sfield{Rationale} \cref{prop:rs:from-signal-to-forecast:scaling}.
\sfield{Horizon and half-life} The composite's; over $h$ periods, the sum of its single-period ICs.
\sfield{Normalisation} Already in return units; checked each month by the Mincer--Zarnowitz slope.
\sfield{Failure modes} An IC estimated in sample; volatilities from a stale risk model; tails extrapolated from a linear fit.
\sfield{Sources} Grinold (1989); \texttt{firm.forecast}; \texttt{rs\_forecast.calibration}.
\end{predictorcard}

\section{Tutorial: why the IR is not what the law promised}\label{tut:rs:from-signal-to-forecast:tut}

\begin{tutorial}
\textbf{Goal.} Calibrate a planted score into expected returns three ways, check the law with Gaussian residuals, then measure the shortfall on the
market's residuals and under a long-only constraint. \textbf{End state:}
\cref{fig:rs:from-signal-to-forecast:calibration,fig:rs:from-signal-to-forecast:shortfall}; the tracking-error table.
\begin{tutsteps}
\item \textbf{\omterm{def:rs:from-signal-to-forecast:curve}{Isotonic regression}} by pooling adjacent violators.
\omcode{code/firm/forecast/firm_forecast.py}{42}{62}{Pool adjacent violators.}\label{lst:rs:from-signal-to-forecast:pav}
\item \textbf{The law and its corrections.}
\omcode{code/firm/forecast/firm_forecast.py}{73}{91}{Effective breadth, transfer coefficient, and the laws.}\label{lst:rs:from-signal-to-forecast:law}
\item \textbf{Run} \texttt{calibration()}, \texttt{law()}, \texttt{tracking\_error\_grid()} and \texttt{fig\_forecast.py}.
\end{tutsteps}
\textbf{What to change next.} Make the IC depend on the name's volatility (higher for small, volatile names) and see the scaling rule miscalibrate;
add a sector-neutrality constraint and measure its \omterm{def:rs:from-signal-to-forecast:breadth}{transfer coefficient}.
\end{tutorial}

\section{Build: the forecast toolkit}\label{bld:rs:from-signal-to-forecast:forecast}

\begin{build}
\bfield{Purpose} Every composite that reaches a portfolio arrives as an expected return, with its calibration checked and its expected information ratio
decomposed.
\bfield{Interface} \texttt{scale\_rule(z, ic, vol)}, \texttt{binned}, \texttt{isotonic}, \texttt{mincer\_zarnowitz}, \texttt{effective\_breadth},
\texttt{transfer\_coefficient(w, alpha, vol)}, \texttt{law\_ir(ic, breadth, tc)}, \texttt{qian\_hua\_ir}, \texttt{ding\_martin\_ir},
\texttt{information\_ratio}.
\bfield{Rules} ICs used for scaling are estimated before the dates they scale; forecasts are in return units per horizon; every backtest reports the
law's promise beside its realised information ratio.
\bfield{Acceptance tests} \texttt{code/firm/forecast/tests/}: the scaling rule and bins by hand; \omterm{def:rs:from-signal-to-forecast:curve}{isotonic regression} on a hand sequence, on unsorted
inputs and on a noisy line; Mincer--Zarnowitz on a planted slope; \omterm{def:rs:from-signal-to-forecast:breadth}{effective breadth} at the extremes; the laws by hand; a \omterm{def:rs:from-signal-to-forecast:breadth}{transfer coefficient} of 1
for the unconstrained book and below 1 when shorts are removed.
\bfield{Stretch} Calibration by regime; IC by volatility bucket; the Ding and Martin form with an estimated conditional IC.
\end{build}

\begin{omsources}
\item R.~C.~Grinold, ``The fundamental law of active management'', \emph{Journal of Portfolio Management} 15(3), 1989.
\item R.~Clarke, H.~de Silva and S.~Thorley, ``Portfolio constraints and the fundamental law of active management'', \emph{Financial Analysts Journal}
58(5), 2002.
\item F.~Zhang, X.~Wang and H.~Cao, ``Turnover-adjusted information ratio'', arXiv:2105.10306, 2021 (the law and its Qian--Hua and Ding--Martin forms).
\item R.~O.~Michaud, D.~N.~Esch and R.~O.~Michaud, ``Estimation error and the fundamental law of active management'', \emph{Journal of Investing}
29(4), 2020.
\item M.~Ayer, H.~D.~Brunk, G.~M.~Ewing, W.~T.~Reid and E.~Silverman, ``An empirical distribution function for sampling with incomplete information'',
\emph{Annals of Mathematical Statistics} 26(4), 1955.
\end{omsources}

\section{Exercises}

\begin{exercise}[$\star$]\label{exo:rs:from-signal-to-forecast:1}
A score of $-1.5$, an IC of 0.04 and a monthly residual volatility of 10\%: what is the expected \omterm{def:rs:anatomy-of-a-predictor:residual}{residual return} for the month? And with a score of
$+2$ and a volatility of 5\%?
\end{exercise}

\begin{exercise}[$\star$]\label{exo:rs:from-signal-to-forecast:2}
What information ratio does the law promise for an IC of 0.03 on 500 names rebalanced monthly? On 50 names rebalanced weekly (52 times a year)?
\end{exercise}

\begin{exercise}[$\star$]\label{exo:rs:from-signal-to-forecast:3}
Fit \omterm{def:rs:from-signal-to-forecast:curve}{isotonic regression} by hand to the values 2, 1, 3, 3, 2, 4 at increasing scores.
\end{exercise}

\begin{exercise}[$\star\star$]\label{exo:rs:from-signal-to-forecast:4}
One hundred bets with an average pairwise correlation of 0.05: what is their \omterm{def:rs:from-signal-to-forecast:breadth}{effective breadth}? What correlation would halve it?
\end{exercise}

\begin{exercise}[$\star\star$]\label{exo:rs:from-signal-to-forecast:5}
A monthly IC has mean 0.04 and standard deviation 0.05. What annual information ratio do Qian and Hua's form and the law with 704 names predict, and
why do they differ?
\end{exercise}

\begin{exercise}[$\star\star$]\label{exo:rs:from-signal-to-forecast:6}
A signal's single-day ICs are 0.02, 0.01 and 0 after the second day, with a daily residual volatility of 2\%. What expected return does a score of $+1$
imply over one day and over a month?
\end{exercise}

\begin{exercise}[$\star\star\star$]\label{exo:rs:from-signal-to-forecast:7}
\emph{Coding.} With \texttt{rs\_forecast.portfolio}, measure the long-only book's \omterm{def:rs:from-signal-to-forecast:breadth}{transfer coefficient} and information ratio at a requested tracking
error of 1\% a month when the benchmark is concentrated (weights proportional to a lognormal size). Why does it change?
\end{exercise}

\begin{exercise}[$\star\star\star$]\label{exo:rs:from-signal-to-forecast:8}
\emph{Find the flaw.} ``Our IC is 0.03 on 3\,000 stocks rebalanced daily, so our information ratio should be $0.03\sqrt{252 \times 3000} = 26$.''
\end{exercise}

\section{Problem: Why Is Our Information Ratio Two-Thirds of the Law's?}

\begin{problem}[{Weekend problem --- a shortfall, decomposed}]\label{pb:rs:from-signal-to-forecast:1}
The chapter's planted forecast: IC 0.05, 704 names, ten years of months, on \texttt{firm.synthmkt}.

\medskip\noindent\textbf{Part I --- Calibration.}
\begin{enumerate}
\item What expected return does a score of $+2$ imply for a name with the median volatility?
\item What IC is estimated on the first five years?
\item Give the Mincer--Zarnowitz slope and $R^2$ of the scaled forecast, and the slope on the raw score.
\item Why is the $R^2$ so small, and why does that not matter?
\item When would you prefer the isotonic calibration to the linear one?
\end{enumerate}

\medskip\noindent\textbf{Part II --- The law.}
\begin{enumerate}[resume]
\item What does the law promise?
\item What does the book realise with Gaussian residuals, and is it consistent with the law?
\item What are the IC's mean and standard deviation with the market's residuals and with Gaussian ones?
\item What \omterm{def:rs:from-signal-to-forecast:breadth}{effective breadth} do they imply?
\end{enumerate}

\medskip\noindent\textbf{Part III --- The constraint.}
\begin{enumerate}[resume]
\item What \omterm{def:rs:from-signal-to-forecast:breadth}{transfer coefficient} does the long-only book reach at a requested tracking error of 4\% a month, and what share of holdings is clipped?
\item What does the generalised law predict, and what is realised?
\item How do the \omterm{def:rs:from-signal-to-forecast:breadth}{transfer coefficient} and information ratio change with the requested tracking error?
\item State the \emph{named result}: the decomposition of the shortfall into \omterm{def:rs:from-signal-to-forecast:breadth}{effective breadth}, \omterm{def:rs:from-signal-to-forecast:breadth}{transfer coefficient} and the noise of the constraint.
\end{enumerate}

\medskip\noindent\textbf{Part IV --- Beyond.}
\begin{enumerate}[resume]
\item What would costs do to the last \omterm{def:rs:market-data-for-research:bar}{bar}?
\item How would a cap-weighted benchmark change the \omterm{def:rs:from-signal-to-forecast:breadth}{transfer coefficient}?
\item Why can adding names reduce the information ratio in practice?
\item What should the firm report for every strategy?
\item Which of the three steps can research improve, and which only portfolio construction?
\item What is the danger of using the law to size a new strategy?
\item In one sentence: what does the law measure?
\end{enumerate}
\end{problem}

\section{Interview questions}

\begin{interviewq}[$\star$ \iqroles{researcher}]\label{iq:rs:from-signal-to-forecast:1}
State the \omterm{def:rs:from-signal-to-forecast:law}{fundamental law of active management} and its assumptions.
\end{interviewq}

\begin{interviewq}[$\star\star$ \iqroles{researcher}]\label{iq:rs:from-signal-to-forecast:2}
How do you turn an alpha score into an expected return?
\end{interviewq}

\begin{interviewq}[$\star\star$ \iqroles{researcher, trader}]\label{iq:rs:from-signal-to-forecast:3}
Your backtest's information ratio is half what the law predicts. Where do you look?
\end{interviewq}

\begin{interviewq}[$\star\star$ \iqroles{researcher}]\label{iq:rs:from-signal-to-forecast:4}
What is the \omterm{def:rs:from-signal-to-forecast:breadth}{transfer coefficient}, and what lowers it?
\end{interviewq}

\begin{interviewq}[$\star\star$ \iqroles{researcher, mle}]\label{iq:rs:from-signal-to-forecast:5}
How would you check that a model's return forecasts are calibrated?
\end{interviewq}

\begin{interviewq}[$\star\star\star$ \iqroles{researcher}]\label{iq:rs:from-signal-to-forecast:6}
Derive the generalised law $\mathrm{IR} = \mathrm{TC}\cdot\mathrm{IC}\sqrt{\mathrm{BR}}$ and say which assumption fails first in practice.
\end{interviewq}
